% !TeX program = lualatex
% halo:begin
% {
%   "version": 1,
%   "publish": true,
%   "course": "mathematical-analysis-i",
%   "section": "3.4",
%   "title": "数学分析（一）3.4：无穷小量与无穷大量",
%   "categories": ["study-notes"],
%   "tags": ["mathematics"],
%   "excerpt": "无穷小量的运算、阶的比较与等价代换，无穷大量与倒数关系，以及曲线渐近线的判定。"
% }
% halo:end
% 来源：IMG_2460—IMG_2463、IMG_2473—IMG_2479.HEIC。
% 教材核对：华东师大《数学分析》第5版上册，PDF第67—73页（书页56—62）。
% 两份根式例题分别保留1+2x²与1−2x²；后一式的极限修正为−4。
\RequirePackage{iftex}
\RequireLuaTeX
\documentclass[UTF8,fontset=fandol,a4paper,12pt]{ctexart}
\usepackage[margin=2.5cm]{geometry}
\usepackage{amsmath,amssymb,amsthm}
\usepackage{enumitem,fancyhdr}
\usepackage[hidelinks]{hyperref}
\newcommand{\R}{\mathbb R}
\newcommand{\U}{U}
\theoremstyle{definition}
\newtheorem*{definition}{定义}
\newtheorem*{theorem}{定理}
\newtheorem*{example}{例题}
\renewcommand{\proofname}{证明}
\setlist[enumerate]{leftmargin=2em,itemsep=0.3em,topsep=0.3em}
\setlength{\headheight}{16pt}
\pagestyle{fancy}
\fancyhf{}
\fancyhead[L]{\small 数学分析（一）}
\fancyhead[R]{\small 3.4\quad 无穷小量与无穷大量}
\fancyfoot[C]{\thepage}
\renewcommand{\headrulewidth}{0.3pt}
\raggedbottom
\title{无穷小量与无穷大量}
\author{Yuchen Zhang}
\date{整理日期：2026年10月22日}
\makeatletter
\renewcommand{\maketitle}{\begin{center}
{\LARGE\bfseries \@title\par}\vspace{0.6em}
\@author\par\vspace{0.2em}{\small\@date\par}
\end{center}}
\makeatother
\begin{document}
\maketitle

\section{无穷小量及其运算}
有限极限 $f(x)\to A$ 可以改写为 $f(x)-A\to0$。把这部分趋于零的误差单独记出，就得到无穷小量的概念。

\begin{definition}
在给定的趋近过程中，若 $f(x)\to0$，则称 $f$ 为该过程中的无穷小量。若 $f$ 在相应的某个去心邻域或远端区间内有界，则称它为有界量。
\end{definition}

无穷小量是具有指定极限性质的函数，不是一个固定的“小数”。趋近过程也是定义的一部分：$x$ 在 $x\to0$ 时为无穷小量，在 $x\to+\infty$ 时则不是。零函数是任一趋近过程中的无穷小量。

例如 $x\to0$ 时，$\sin x$、$e^x-1$、$\ln(1+x)$、$x^2$、$1-\cos x$ 都是无穷小量；$\sqrt{1-x}$ 则是 $x\to1^-$ 时的无穷小量。一般地，
\[
 f(x)\to A\quad\Longleftrightarrow\quad
 f(x)=A+\alpha(x),\qquad\alpha(x)\to0.
\]

在同一趋近过程中，有限个无穷小量的和、差、积仍是无穷小量。这些结论由四则运算法则得到。无穷小量与有界量相乘也仍是无穷小量，即使有界因子本身没有极限。

\begin{proof}
设 $\alpha(x)\to0$，且在相应的某个邻近区域内 $|\beta(x)|\le M$，其中取 $M>0$。任取 $\varepsilon>0$，由 $\alpha(x)\to0$，可进一步限制自变量，使 $|\alpha(x)|<\varepsilon/M$。于是
\[
 |\alpha(x)\beta(x)|\le M|\alpha(x)|<\varepsilon.
\]
故 $\alpha(x)\beta(x)\to0$。
\end{proof}

\begin{example}
比较 $x\sin(1/x)$ 在 $x\to0$ 和 $x\to\infty$ 时的极限。
\end{example}
\begin{proof}[解]
当 $x\to0$ 时，$x$ 为无穷小量，$\sin(1/x)$ 为有界量，所以
\[
 \lim_{x\to0}x\sin\frac1x=0.
\]
当 $x\to\infty$ 时，$x$ 不再是无穷小量，不能套用刚才的性质。此时改写为
\[
 x\sin\frac1x=\frac{\sin(1/x)}{1/x}\to1.
\]
同一个表达式在不同的趋近过程中可以有不同的极限。
\end{proof}

\section{无穷小量阶的比较}
两个量都趋于零，并不表示它们的大小相近。比较它们的商，可以进一步区分趋零的速度。以下固定同一趋近过程，并在所写比值出现时假定分母最终非零。

\subsection*{高阶无穷小与 $o$、$O$ 记号}
若
\[
 \frac{f(x)}{g(x)}\to0,
\]
记为 $f(x)=o(g(x))$。当 $f,g$ 都是无穷小量时，称 $f$ 为 $g$ 的高阶无穷小量，或称 $g$ 为 $f$ 的低阶无穷小量。例如
\[
 x^2=o(x),\qquad 1-\cos x=o(x)\qquad(x\to0).
\]
第二式可由 $(1-\cos x)/x=[(1-\cos x)/x^2]x\to0$ 得到。

若存在常数 $C>0$，使趋近过程中最终有
\[
 |f(x)|\le C|g(x)|,
\]
则记为 $f(x)=O(g(x))$。它只要求商有界，不要求商有极限。例如 $x^2=O(x)$、$1-\cos x=O(x^2)$（$x\to0$）。特别地，$f=o(1)$ 表示 $f\to0$，$f=O(1)$ 表示 $f$ 在该趋近过程中局部有界。

这里的 $o(g)$、$O(g)$ 分别表示一类函数，等号表示 $f$ 属于相应的函数类。由定义，$f=o(g)$ 必有 $f=O(g)$，反之不一定成立。

\subsection*{同阶与等价}
\begin{definition}
设 $f,g$ 都是无穷小量。若存在常数 $K,L>0$，使趋近过程中最终有
\[
 K\le\left|\frac{f(x)}{g(x)}\right|\le L,
\]
则称 $f,g$ 为同阶无穷小量。若进一步有
\[
 \frac{f(x)}{g(x)}\to1,
\]
则称 $f,g$ 为等价无穷小量，记为 $f(x)\sim g(x)$。
\end{definition}

若 $f/g\to c\ne0$，则 $f,g$ 必为同阶无穷小量，但同阶不要求这个商有极限。例如当 $x\to0$ 时，
\[
 f(x)=x\left(2+\sin\frac1x\right),\qquad g(x)=x
\]
都趋于零，而 $1\le|f(x)/g(x)|\le3$，所以它们同阶；商在 $1$ 与 $3$ 之间振荡，没有极限，因而不等价。

等价比同阶的要求更强，并且可以改写为
\[
 f\sim g\quad\Longleftrightarrow\quad f=g+o(g),
\]
因为 $(f-g)/g=f/g-1$。这说明用 $g$ 代替 $f$ 时，省略的差相对于 $g$ 是高阶无穷小。

\subsection*{常用等价无穷小}
前一篇的比值极限现在可以汇成一组等价关系。当 $x\to0$ 时，
\[
 \begin{gathered}
 \sin x\sim x,\quad \tan x\sim x,\quad
 \arcsin x\sim x,\quad \arctan x\sim x,\\
 e^x-1\sim x,\qquad\ln(1+x)\sim x,\\
 1-\cos x\sim\frac{x^2}{2},\qquad
 (1+x)^\alpha-1\sim\alpha x\quad(\alpha\ne0).
 \end{gathered}
\]
最后一式要求 $\alpha\ne0$，因为 $\alpha=0$ 时两边恒为零，不能形成商趋于 $1$ 的等价关系；不过极限公式 $[(1+x)^\alpha-1]/x\to\alpha$ 对所有实数 $\alpha$ 都成立。

\section{等价代换与和差式}
等价关系能简化乘积和商的极限，因为代换前后的比值趋于 $1$。

\begin{theorem}[等价代换]
设 $f\sim\widetilde f$、$g\sim\widetilde g$，且各比值最终有定义。若 $\widetilde f/\widetilde g\to A\in\R$，则
\[
 \frac{f}{g}\to A.
\]
反之亦然。等价因子在乘积中的代换也有同样的性质。
\end{theorem}
\begin{proof}
由
\[
 \frac{f}{g}
 =\frac{f}{\widetilde f}\,
   \frac{\widetilde f}{\widetilde g}\,
   \frac{\widetilde g}{g},
\]
右端三个因子分别趋于 $1,A,1$，故原商趋于 $A$。反向只需交换原量与等价量。乘积情形同理。
\end{proof}

\begin{example}
计算 $\displaystyle\lim_{x\to0}\frac{\arctan x}{\sin4x}$。
\end{example}
\begin{proof}[解]
由 $\arctan x\sim x$、$\sin4x\sim4x$，得
\[
 \lim_{x\to0}\frac{\arctan x}{\sin4x}
 =\lim_{x\to0}\frac{x}{4x}=\frac14.
\]
\end{proof}

两个等价无穷小量相减时，原来省略的部分可能成为主要部分，因此不能在和差内部任意逐项代换。

\begin{example}
计算 $\displaystyle\lim_{x\to0}\frac{\tan x-\sin x}{\sin(x^3)}$。
\end{example}
\begin{proof}[解]
若把分子中的 $\tan x$、$\sin x$ 都直接换为 $x$，便会把分子变成 $x-x=0$，丢掉两者之差的信息。正确做法是先用恒等式
\[
 \tan x-\sin x=\frac{\sin x}{\cos x}(1-\cos x)
\]
把差化为乘积。再用 $\sin x\sim x$、$1-\cos x\sim x^2/2$、$\sin(x^3)\sim x^3$，得到
\[
 \lim_{x\to0}\frac{\tan x-\sin x}{\sin(x^3)}
 =\lim_{x\to0}\frac1{\cos x}\frac{x\cdot(x^2/2)}{x^3}
 =\frac12.
\]
\end{proof}

这并不妨碍我们在各项极限都存在时使用和差法则。若整个待求表达式能拆成几个各有有限极限的项，就分别计算，再合并结果。

\begin{example}
计算
\[
 \lim_{x\to0}\frac{\tan x-x^2}{\sin x},\qquad
 \lim_{x\to0}\frac{\tan x\pm(1-\cos x)}{\sin x}.
\]
\end{example}
\begin{proof}[解]
第一式拆成
\[
 \frac{\tan x}{\sin x}-\frac{x^2}{\sin x}
 =\frac1{\cos x}-x\frac{x}{\sin x}\to1-0=1.
\]
第二式中，
\[
 \frac{1-\cos x}{\sin x}
 =\frac{1-\cos x}{x^2}\,\frac{x}{\sin x}\,x\to0,
\]
而 $\tan x/\sin x=1/\cos x\to1$，所以两种符号下的极限均为 $1$。
\end{proof}

\begin{example}
分别计算
\[
 \lim_{x\to0}\frac{\sqrt[3]{1+2x^2}-1}
 {\arcsin(x/2)\,\arctan(x/3)},\qquad
 \lim_{x\to0}\frac{\sqrt[3]{1-2x^2}-1}
 {\arcsin(x/2)\,\arctan(x/3)}.
\]
\end{example}
\begin{proof}[解]
前一篇得到的幂函数增量极限，也可以写成等价关系，用来处理根式。由 $(1+u)^{1/3}-1\sim u/3$，有
\[
 \sqrt[3]{1+2x^2}-1\sim\frac23x^2,\qquad
 \sqrt[3]{1-2x^2}-1\sim-\frac23x^2.
\]
两式的分母均满足
\[
 \arcsin(x/2)\,\arctan(x/3)\sim\frac x2\frac x3=\frac{x^2}{6}.
\]
故两个极限依次为 $4$、$-4$。根号内的正负号决定分子主要部分的符号，不能省去。
\end{proof}

\section{无穷大量}
前面讨论的是趋于零的量。还有一类函数，其函数值的绝对值会随着自变量的趋近而超过任意给定的正数，这就引出无穷大量的概念。

\begin{definition}
设 $f$ 在 $x_0$ 的某个去心邻域内有定义。若对任意 $G>0$，存在 $\delta>0$，使
\[
 0<|x-x_0|<\delta\quad\Longrightarrow\quad |f(x)|>G,
\]
则称 $f$ 为 $x\to x_0$ 时的无穷大量，记为 $f(x)\to\infty$。
若把结论分别改为 $f(x)>G$ 或 $f(x)<-G$，则分别记为 $f(x)\to+\infty$ 或 $f(x)\to-\infty$。
\end{definition}

这里不带符号的 $\infty$ 表示绝对值趋于正无穷，不要求符号固定。例如
\[
 \lim_{x\to0}\frac1x=\infty,\qquad
 \lim_{x\to0^+}\frac1x=+\infty,\qquad
 \lim_{x\to0^-}\frac1x=-\infty.
\]
第一个式子不表示存在有限实数极限，也不表示双侧极限为 $+\infty$。无穷远处和单侧的定义，只需替换相应的趋近区域。

由定义还得到倒数关系：若 $f$ 最终非零，则
\[
 f(x)\to\infty\quad\Longleftrightarrow\quad\frac1{f(x)}\to0.
\]
因为 $|f(x)|>G$ 与 $|1/f(x)|<1/G$ 等价。

无界比无穷大的要求弱。例如 $x\sin x$ 在任意远的右半区间上都无界，但在 $x=2n\pi$ 时函数值为零，所以它在 $x\to+\infty$ 时不是无穷大量。

\section{曲线的渐近线}
函数极限还可以描述图像怎样接近一条直线。有限极限可以给出水平渐近线，无穷极限可以给出垂直渐近线；图像在无穷远处还可能接近一条斜直线。

\subsection*{水平渐近线与垂直渐近线}
若 $f(x)\to A\in\R$（$x\to+\infty$ 或 $x\to-\infty$），则直线 $y=A$ 是相应方向的水平渐近线。两个方向应分别考察，所得直线可以不同。

若 $x\to x_0^+$ 或 $x\to x_0^-$ 时 $|f(x)|\to+\infty$，则 $x=x_0$ 是垂直渐近线；只要一侧满足即可。几何上，点 $(x,f(x))$ 到原点的距离趋于无穷，而它到直线 $x=x_0$ 的距离 $|x-x_0|$ 趋于零。

\subsection*{斜渐近线}
若存在有限常数 $k\ne0$、$b$，使 $x\to+\infty$ 时
\[
 f(x)-(kx+b)\to0,
\]
则 $y=kx+b$ 是该方向的斜渐近线。点 $(x,f(x))$ 到直线的距离为
\[
 \frac{|f(x)-(kx+b)|}{\sqrt{1+k^2}},
\]
所以这一极限正好表示点到直线的距离趋于零。$x\to-\infty$ 的情形同理，仍应单独考察。

为了确定 $k,b$，先由 $f(x)-kx\to b$ 得
\[
 \frac{f(x)}x=k+\frac{f(x)-kx}{x}\to k.
\]
因此可以依次计算
\[
 \boxed{k=\lim_{x\to+\infty}\frac{f(x)}x,\qquad
 b=\lim_{x\to+\infty}[f(x)-kx].}
\]
反过来，如果这两个极限存在有限值且 $k\ne0$，则 $f(x)-(kx+b)\to0$，从而确有该斜渐近线。只算出斜率 $k$ 还不够，必须继续检查截距 $b$ 的极限。

\begin{example}
求 $\displaystyle f(x)=\frac{x^3}{x^2+2x-3}$ 的渐近线。
\end{example}
\begin{proof}[解]
先分解分母：$x^2+2x-3=(x+3)(x-1)$。分母仅在 $x=-3,1$ 处为零，且两点处分子均不为零。考察各侧的符号可得
\[
 \begin{aligned}
 &\lim_{x\to-3^-}f(x)=-\infty,\qquad
 \lim_{x\to-3^+}f(x)=+\infty,\\
 &\lim_{x\to1^-}f(x)=-\infty,\qquad
 \lim_{x\to1^+}f(x)=+\infty.
 \end{aligned}
\]
故垂直渐近线为 $x=-3$、$x=1$。其他有限点的分母非零，函数有有限极限，不会产生垂直渐近线。

在 $x\to+\infty$ 与 $x\to-\infty$ 两个方向，均有
\[
 \frac{f(x)}x=\frac{x^2}{x^2+2x-3}\to1,
\]
\[
 f(x)-x=\frac{-2x^2+3x}{x^2+2x-3}\to-2.
\]
所以两端的斜渐近线都是 $y=x-2$。也可用多项式除法直接核对：
\[
 f(x)=x-2+\frac{7x-6}{x^2+2x-3},
\]
最后一项趋于零。因而 $x\to+\infty$ 时 $f(x)\to+\infty$，$x\to-\infty$ 时 $f(x)\to-\infty$，没有水平渐近线。
\end{proof}

\end{document}
